\sscp{\tsc{East Bomberai}}{11-05-02}
The \tsc{East Bomberai} node of \tsc{Tanimbar-Bomberai} contains
Irarutu, Kuri (a.k.a Nabi), Arguni, Bedoanas (a.k.a.
Yarik), and Erokwanas (a.k.a. Kambran, Damberang-Goras).
The \tsc{East Bomberai} languages have had extensive contact
with neighbouring Papuan languages,
in particular Iha and Mbaham.
This contact has resulted in borrowing in both directions
\citep{UsherSchapper2018,UsherSchapper2022}.
\citet{UsherSchapper2018} focus on borrowing from the Austronesian
languages into the Papuan languages,
but borrowings in the other direction also occur.
In particular, \citet{UsherSchapper2022} include the highly significant
Proto-Greater West Bomberai reconstructions *mandəl ‘bat’
and *maŋgʷVl ‘banana’ discussed in
\srf{14-04-04} and \srf{02-03-01-07} respectively.
Two examples which are more locally restricted are Iha \it{hɛir},
Mbaham \it{sey(i)r} > Arguni \it{sair},
Bedoanas, Erokwanas \it{sai} all ‘fish’,
as well as Iha \it{tɔqar},
Mbaham \it{tɔːqar} both ‘bone’ > Arguni \it{tór},
Bedoanas \it{tóʔor}, Erokwanas \it{toʔar} all ‘thorn,
bone’, and Irarutu \it{tor} [torə] ‘bone’.\footnote{
		Arguni \it{tór},
		Bedoanas \it{tóʔor}, Erokwanas \it{toʔar} all ‘thorn,
		bone’, and Irarutu \it{tor} [torə] ‘bone’ might alternately be
		from PMP *tuqəlan ‘condylar bone; bone of fauna exclusive of
		fish’, though we would expect *q > {\0} (not *q > \it{ʔ}) in Bedoanas
		and Erokwanas.} 
The heavily Papuanised nature of these languages,
combined with scarcity of data until recent years,
has meant that the position of these languages has remained unclear.
\citet[272]{Blust1993} placed Arguni with the \tsc{Bomberai Tip}
languages (and thus within his CMP) noting that it “[{\ldots}]
is the most divergent member of the group.” \citet{Blust1993}
left Irarutu unclassified within CEMP,
and had no access to data from Bedoanas or Erokwanas to make any assessment.
\citet{Kamholz2014}, in his investigation of the SHWNG subgroup,
also considered the position of these languages focusing on Irarutu.
He came to no firm conclusion but excluded them from his examination of SHWNG.

\largerpage
Evidence that these languages do not belong within SHWNG comes
from the reflexes of *j,
as exemplified in \trf{11-04} and \trf{11-22}.
SHWNG languages show a merger of *s/*j \citep[52f]{Kamholz2014}.
Unlike SHWNG languages, *j
and *s do not usually merge in \tsc{East Bomberai}.
The reflexes of *j are discussed in more detail in the following sections.
Evidence that these languages belong within STB comes from shared
merger of *z/*d > **d (\srf{11-01-01}).

\begin{table}\tcp{\tsc{East Bomberai} *l > \it{r}}{11-30}
%\begin{adjustbox}{width=\textwidth}
	\begin{threeparttable}\st{6pt}
	\begin{tabular}{llllll}\lsptoprule
*gloss	&	sky	&	head	&	ten	&	root	&	road, path\\
PMP	&	*\tbr{l}aŋit	&	*qu\tbr{l}uh	&	*saŋapu\tbr{l}uq	&	*da\tbr{l}ij\su{Δ}	&	*za\tbr{l}an\\	\midrule
Arguni	&	{\hr}\tbr{r}ágit	&	{\hr}á\tbr{r}u-g	&	{\hr}sambú\tbr{r}e	&	{\hr}da\tbr{r}e(g)	&	{\hr}rá\tbr{r}in\\
Bedoanas	&	{\hr}\tbr{r}aʔit	&	{\hr}i\tbr{r}i atu\su{†}	&	{\hr}sambú\tbr{r}a	&	{\hr}da\tbr{r}a, nda\tbr{r}ao	&	{\hr}rá\tbr{r}in\\
Erokw.	&	{\hr}\tbr{r}áʔit	&	{\hr}a\tbr{r}u puat\su{†}	&	{\hr}sambú\tbr{r}a	&	{\hr}dá\tbr{r}a	&	{\hr}ró\tbr{r}in\\
Irarutu	&	{\hr}\tbr{r}agt	&	{\hr}\tbr{r}ü \su{‡}	&	\cg{(fradru)}	&	\cg{(war)}	&	{\hr}ra\tbr{r}n\\
Kuri	&		&	\cg{(tə̆bon)}	&	\cg{(βradu)}	&	\cg{(war)}	&	{\hr}ra\tbr{r}n\\
		\lspbottomrule
	\end{tabular}
		\begin{tablenotes}\tnsize
			\item[†]	{Bedoanas \it{atu}
								and Erokwanas \it{puat} are probably from *batu ‘stone’.}
			\item[‡]	{Irarutu \it{rü} is glossed ‘source’ \citep[222]{Jackson2014}.
								The vowel is a front rounded vowel.}
			\item[Δ]	{Reflexes of *dalij have *d = \it{d} against normal *d > \it{r}.
								However, compare also the reflexes of *z (\trf{11-03}) for which one two forms
								have *z > \it{r} and one has *z > \it{d}.
								Arguni \it{dareg} is from \cite{Narfafan2011} and \it{dare}
								from \citet[215]{SmitsVoorhoeve1992a}.}
		\end{tablenotes}
	\end{threeparttable}
%\end{adjustbox}
\end{table}

The main defining sound change of \tsc{East Bomberai} is shared
*l > \it{r}.\footnote{
		The Northeast Arguni dialect of Irarutu
		has subsequent **r > \it{l} \citep[8]{Jackson2014}.}
This change also occurs in neighbouring Kowiai (\srf{11-03-02}).
Examples are given in \trf{11-25},
with an additional four examples given later in \trf{11-30}.
To these we can add a number of words for which only the Arguni
reflex is currently known: *kapa\tbr{l} > \it{maʔápe\tbr{r}}
‘thick’ (with irregular *p = \it{p}),
*qa\tbr{l}ipan > \it{a\tbr{r}ifen} ‘centipede’,
and *tə\tbr{l}ən > \it{ute\tbr{r}en} ‘swallow’.

\begin{table}\tcp{\tsc{East Bomberai} *ŋ > \it{g}}{11-31}
\begin{adjustbox}{width=\textwidth}
	\begin{threeparttable}\st{2pt}
	\begin{tabular}{lllllll}\lsptoprule
*gloss	&	name	&	cry	&	ear	&	finger, toe	&	branch	&	sky\\
PMP	&	*\tbr{ŋ}ajan	&	*ta\tbr{ŋ}is	&	*ta\tbr{ŋ}ila/*tali\tbr{ŋ}a	&	*ta\tbr{ŋ}an	&	*sa\tbr{ŋ}a	&	*la\tbr{ŋ}it\\	\midrule
Yamdena	&	{\hr}\tbb{ŋ}ar-e	&	{\hr}-tasi\tbb{ŋ}	&	{\hr}tli\tbb{ŋ}a-\tbb{ŋ}w	&	{\hr}ta\tbb{ŋ}a-ny	&	{\hr}sa\tbb{ŋ}a-ny	&	{\hr}la\tbb{ŋ}it\\	\midrule
Fordata	&	{\hr}\tbb{n}ara-\tbb{n}	&	{\hr}\tbb{n}-ta\tbb{n}it	&	\cg{(aru-n)}	&		&	{\hr}a\tbb{n}a-\tbb{n}	&	{\hr}la\tbb{n}it\\
Kei	&	\cg{(mema-n)}	&	{\hr}tna\tbb{n}it	&	\cg{(aru-n)}	&		&	{\hr}ha\tbb{ŋ}a	&	{\hr}la\tbb{n}it\\
Onin	&	{\hr}\tbr{g}ara-r	&	\cg{(a-gaga)}	&	{\hr}tani\tbr{g}a-n	&	{\hr}nemi-n ta\tbr{k}an	&		&	\cg{(laŋit)}\su{†}\\
Sekar	&	{\hr}\tbr{g}ara-r	&	{\hr}a-ta\tbr{g}is	&	{\hr}tani\tbr{g}a-n	&	{\hr}nima-n ta\tbr{g}an	&		&	\cg{(laŋit)}\su{†}\\
Uruang.	&	{\hr}(\tbr{ŋ})\tbr{g}aran 	&	\cg{(a-rekar)}	&	{\hr}tani\tbr{ŋ}a-n	&	{\hr}ta\tbr{ŋ}karek	&	{\hr}ra\tbr{ŋg}as	&	{\hr}na\tbr{g}it\su{†}\\	\midrule
Arguni	&	{\hr}\tbr{g}ádin	&	{\hr}t\<y\>á\tbr{g}īs	&	{\hr}ti\tbr{g}yánáne	&	{\hr}tá\tbr{g}e	&	{\hr}sá\tbr{g}e	&	{\hr}rá\tbr{g}it\\
Bedoanas	&	{\hr}\tbr{g}adiŋ	&	{\hr}t\<i\>á\tbr{ʔ}is	&		&	\cg{(baʔya)}	&		&	{\hr}ra\tbr{ʔ}it\\
Erokw.	&	{\hr}\tbr{g}adin	&	{\hr}t\<i\>á\tbr{ʔ}is	&	{\hr}ti\tbr{ʔ}inyanya	&	\cg{(rumwa)}	&		&	{\hr}rá\tbr{ʔ}it\\
Irarutu	&	\cg{(no)}	&	{\hr}n-ta\tbr{g}	&	{\hr}t\tbr{g}ra	&	\cg{(fra ntü)}	&	\cg{(ran)}	&	{\hr}ra\tbr{g}t\\
Kuri	&	\cg{(inɔi)}	&	{\hr}itə-ta\tbr{g}	&	{\hr}tə̆\tbr{g}ra	&		&	{\hr}sa\tbr{g}a	&	\\
		\lspbottomrule
	\end{tabular}
		\begin{tablenotes}\tnsize
			\item[†] {The Onin and Sekar words for ‘sky’ are Malay borrowings.
								Uruangnirin \it{nagit} ‘sky’ is from \citet[22]{SmitsVoorhoeve1992b}.
								\citet{Visser2023} has \it{laŋit}.}
		\end{tablenotes}
	\end{threeparttable}
\end{adjustbox}
\end{table}

Another consonant change shared between the \tsc{East Bomberai}
languages is *ŋ > \it{g} with subsequent *g > \it{ʔ} in Bedoanas and Erokwanas.
Examples are given in \trf{11-31}.
This change is not an exclusively shared innovation.
It also occurs in \tsc{Bomberai Tip} languages.
Onin and Sekar have *ŋ > \it{g}
and Uruangnirin has *ŋ > \it{ŋg, g, ŋ}.
Kowiai also has *g > \it{ŋg, g} (\srf{11-03-02}),
Teor has *ŋ > \it{g} (\srf{11-04}),
and \tsc{East Seram} has similar medial *ŋ
> \it{k} (\srf{11-02}).

\tsc{East Bomberai} languages usually share word-initial *b >
**p, with subsequent **p > {\0} in Bedoanas.
Irarutu has subsequent **b > \it{f} thus indicating a merger of PMP *b/*p.
Kuri also appears to merge *b/*p with a variety of realisations
(probably allophonic) including: [b],
[β], [f], and [v].
Whether the merger in Kuri can be reconciled with *b > **p posited
for other languages is unclear.
Examples of PMP *b in \tsc{East Bomberai} are given in \trf{11-32}.
Onin and Uruangnirin of \tsc{Bomberai Tip} also have *b > \it{p.}

\begin{table}
	\centering\tcp{\tsc{East Bomberai} *b}{11-32}
	\st{6pt}\begin{tabular}{llllll}\lsptoprule
*gloss	&	moon	&	body hair, feather	&	new	&	stone	&	bow\\
PMP	&	*\tbr{b}ulan	&	*\tbr{b}ulu	&	*\tbr{b}aqəRu	&	*\tbr{b}atu	&	*\tbr{b}usuR\\	\midrule
Arguni	&	{\hr}\tbr{p}úrin	&	{\hr}pu-\tbr{p}ure	&	{\hr}pu-\tbr{p}uaru	&	{\hr}\tbr{p}uát	&	{\hr}u\tbr{p}úsir\\
Bedoanas	&	{\hr}úrin	&	{\hr}ur-ura	&	{\hr}ur-uri	&	{\hr}uat	&	\\
Erokw.	&	{\hr}\tbr{p}úrin	&	{\hr}pu-\tbr{p}ura	&	{\hr}\tbr{p}uari	&	{\hr}\tbr{p}uat	&	\\
Irarutu	&	\cg{(seba)}	&	{\hr}\tbr{f}ru	&	\cg{(bunat)}	&	\cg{(kami)}	&	\cg{(fae)}\\
Kuri	&	\cg{(sarabɛ)}	&	{\hr}ˈ\tbr{β}aru, \tbr{b}aru	&	\cg{(aˈbɔb)}	&	\cg{(kɔbaːb)}	&	\cg{(uˈagubɔ)}\\
		\lspbottomrule
	\end{tabular}
\end{table}

In addition to such examples of *b > **p,
there are also at least two examples of *b > \it{(m)b} in Arguni,
Bedoanas, and Erokwanas in our data: *\tbr{b}abuy/*\tbr{b}əRək
‘pig’> Arguni, Bedoanas \it{\tbr{mb}}o, Erokwanas \it{\tbr{b}o},
and *\tbr{b}uaq ‘fruit’ > Arguni \it{\tbr{mb}ú},
Bedoanas, Erokwanas \it{\tbr{mb}u} (Irarutu \it{fu} ‘fruit’).
The reflexes of word medial *b are unclear due to scarcity of examples.
There is one example of *b > \it{p}; *Ri\tbr{b}u ‘thousand’ >
Arguni \it{ri\tbr{p}isye}, Bedoanas \it{ri\tbr{p}isa},
and one example of *b > \it{f}; *tə\tbr{b}uh
> Arguni \it{to\tbr{f}} ‘sugarcane’.

As in many languages of the region reflexes of *ma-p in \tsc{East
Bomberai} are different to those of *b and *p.
Only two items have been identified: *ma-putiq ‘white’ > Irarutu
\it{\tbr{bf}ut} [\tbr{mbəf}ut], Kuri \it{\tbr{maβ}utə̆}, as well as *ma-panas ‘hot,
warm’ > Arguni \it{mba{\tl}\tbr{mb}anes},
Bedoanas, Erokwanas \it{ba{\tl}\tbr{b}anas,} and perhaps Irarutu
\it{\tbr{w}in}, Kuri \it{a\tbr{β}en}.

The reflexes of final vowels also provide evidence for \tsc{East Bomberai}.
Most final syllables are deleted in these languages.
However, they also have epenthesis of final vowels after word
final consonants to various extents.
Thus, loss of final vowels might be best viewed as reduction
rather than complete deletion.
Final syllable deletion (or reduction)
and epenthesis of final vowels are more common in \tsc{Irarutu-Kuri}
than in other \tsc{East Bomberai} languages.
Final vowel epenthesis in Irarutu is described by \citet[653]{Voorhoeve1995a}
as “[{\ldots}] a weak vocalic release.
Usually this is a mid-central vowel [ə],
but it can also be a weak echo-vowel”.
Examples of deletion of final vowels in \tsc{East Bomberai} are given in \trf{11-33}.

\begin{table}\tcp{\tsc{East Bomberai} *V > {\0} /{\gap}{\#}}{11-33}
%\begin{adjustbox}{width=\textwidth}
	\st{2.5pt}\begin{tabular}{llllllll}\lsptoprule
*gloss	&	male	&	die	&	banana	&	faeces	&	sugarcane	&	louse	&	breast\\
PMP	&	*maRuqan\tbr{ay}	&	*mat\tbr{ay}	&	*punt\tbr{i}	&	*taq\tbr{i}	&	*təb\tbr{u}h	&	*kut\tbr{u}	&	*sus\tbr{u}\\	\midrule
Arguni	&	{\hr}máran	&	{\hr}ga-mat	&	{\hr}fúd	&	{\hr}tā	&	{\hr}tof	&	{\hr}út	&	{\hr}sus\\
Bedoanas	&	{\hr}máran	&	{\hr}mú-mat	&	\cg{(maʔo)}	&		&		&	{\hr}ota	&	\cg{(kókon)}\\
Erokw.	&	{\hr}máran	&	{\hr}í-m\<y\>at	&	\cg{(maʔo)}	&	{\hr}ta	&		&	{\hr}óta	&	\cg{(kákon)}\\
Irarutu	&	{\hr}mran	&	{\hr}n-mat	&	{\hr}fud	&	{\hr}ta	&	{\hr}tof	&	{\hr}ut	&	{\hr}sus\\
	&	[məranə]	&	[nəmatə]	&	[fudə]	&	[ta]	&	[tofə]	&	[utə]	&	[susə]\\
Kuri	&	{\hr}bran	&	{\hr}inə-mat	&	\cg{(wɛˈrɔm)}	&		&	\cg{(wɛtit)}	&	\cg{(təni)}	&	{\hr}susɛ\\
		\lspbottomrule
	\end{tabular}
%\end{adjustbox}
\end{table}

Examples of epenthesis of final vowels in \tsc{East Bomberai}
languages (apart from Irarutu
and Kuri) include: *Ratus ‘hundred’ > Arguni \it{rátis\tbr{e}},
Bedoanas \it{rátis\tbr{a}}, Erokwanas \it{ratis\tbr{a}},
*kutu ‘louse’ > **ut > Arguni \it{út},
Bedoanas, Erokwanas \it{ot\tbr{a}}, *sa-ŋa-puluq ‘ten’ > **sambur
> Arguni \it{sambúr\tbr{e}}, Bedoanas, Erokwanas \it{sambúr\tbr{a}},
and *bulu ‘body hair,
feather’ > **pur > Arguni,
\it{pu-pur\tbr{e}} Bedoanas \it{ur-ur\tbr{a}}, Erokwanas,
\it{pu-pur\tbr{a}}.
